% GENERATED FILE. Edit canonical lessons, then run npm run generate:books.
% canonical-source-hash: fnv1a64:11a4f32e5721bf39
% canonical-lessons: RU-C299-nosilki, RU-C299-sirena, RU-C299-trevoga, RU-C299-dezhurnyy, RU-C299-vyzyvat

\chapter{The emergency call}
\label{ch:ru-the-emergency-call}
\hlchaptermodality{299}

\begin{chapteropening}
\textbf{By the end of this chapter:} \emph{I can talk about a stretcher, a siren, an alarm and the person on duty, and say that someone called for help.}

Complete the last lesson of chapter 299: I can talk about a stretcher, a siren, an alarm and the person on duty, and say that someone called for help.
\end{chapteropening}

\section[\texorpdfstring{\ru{носилки} (nosílki)}{носилки (nosílki)}]{\ru{носилки} (nosílki) — a stretcher}

\label{lesson:RU-C299-nosilki}



\begin{quote}
Before the new one: say the Russian for to come out unhurt, then the Russian for to overturn.
\end{quote}

\subsection*{You'll want to know: \ru{носилки}}
\textbf{\ru{носилки}} (\emph{nosílki}) — \textquotedblleft{}a stretcher\textquotedblright{}.

\begin{quote}
\ru{Раненого} \ru{несли} \ru{на} \ru{носилках}.
\end{quote}

\begin{quote}
\emph{Ránenovo neslí na nosílkakh.}
\end{quote}

\begin{quote}
The wounded man was carried on a stretcher.
\end{quote}

Plural only: \emph{\ru{носилки}}, genitive \emph{\ru{носилок}}.

\begin{cousinweb}
\emph{\ru{носилки}} is built on \emph{\ru{носить}} \textquotedblleft{}to carry\textquotedblright{}, which is inherited from Proto-Slavic \emph{nositi}, from PIE \emph{h\textsubscript{1}ne\'{k}-} \textquotedblleft{}to carry\textquotedblright{} (EtymDB).
\end{cousinweb}

\begin{grammarlens}[title={What you've built}]
One more word for telling, explaining and describing at B1.
\end{grammarlens}

\subsection*{Your turn}
\begin{itemize}
\raggedright
  \item \emph{Say it:} \emph{nosílki}
  \item \emph{Say it:} \emph{nosílki}, once more
  \item \emph{Say it:} a sentence of your own with \emph{nosílki}
  \item \emph{Recall:} say the Russian for to come out unhurt, then the Russian for to overturn, then say \emph{nosílki} again
\end{itemize}

\begin{tcolorbox}[breakable,colback=teal!4,colframe=teal!35!black,title={Before you move on}]
What does \ru{носилки} mean? (\textquotedblleft{}A stretcher\textquotedblright{}.) Say it once more.
\end{tcolorbox}
\section[\texorpdfstring{\ru{сирена} (siréna)}{сирена (siréna)}]{\ru{сирена} (siréna) — a siren}

\label{lesson:RU-C299-sirena}



\begin{quote}
Before the new one: say the Russian for to overturn, then the Russian for a stretcher.
\end{quote}

\subsection*{You'll want to know: \ru{сирена}}
\textbf{\ru{сирена}} (\emph{siréna}) — \textquotedblleft{}a siren\textquotedblright{}.

\begin{quote}
\ru{Вдалеке} \ru{выла} \ru{сирена}.
\end{quote}

\begin{quote}
\emph{Vdaleké výla siréna.}
\end{quote}

\begin{quote}
A siren was wailing in the distance.
\end{quote}

Feminine.

\begin{cousinweb}
\emph{\ru{сирена}} is usually traced through French \emph{sirène} to Greek \emph{Seir\'{\={e}}n}, a sea creature of legend whose song drew sailors; EtymDB lists no parent.
\end{cousinweb}

\begin{grammarlens}[title={What you've built}]
One more word for telling, explaining and describing at B1.
\end{grammarlens}

\subsection*{Your turn}
\begin{itemize}
\raggedright
  \item \emph{Say it:} \emph{siréna}
  \item \emph{Say it:} \emph{siréna}, once more
  \item \emph{Say it:} a sentence of your own with \emph{siréna}
  \item \emph{Recall:} say the Russian for to overturn, then the Russian for a stretcher, then say \emph{siréna} again
\end{itemize}

\begin{tcolorbox}[breakable,colback=teal!4,colframe=teal!35!black,title={Before you move on}]
What does \ru{сирена} mean? (\textquotedblleft{}A siren\textquotedblright{}.) Say it once more.
\end{tcolorbox}
\section[\texorpdfstring{\ru{тревога} (trevóga)}{тревога (trevóga)}]{\ru{тревога} (trevóga) — an alarm; anxiety}

\label{lesson:RU-C299-trevoga}



\begin{quote}
Before the new one: say the Russian for a stretcher, then the Russian for a siren.
\end{quote}

\subsection*{You'll want to know: \ru{тревога}}
\textbf{\ru{тревога}} (\emph{trevóga}) — \textquotedblleft{}an alarm; anxiety\textquotedblright{}.

\begin{quote}
\ru{Ночью} \ru{в} \ru{школе} \ru{зазвенел} \ru{сигнал} \ru{тревоги}.
\end{quote}

\begin{quote}
\emph{Nóchyu v shkóle zazvenél signál trevógi.}
\end{quote}

\begin{quote}
At night the alarm went off in the school.
\end{quote}

Also \textquotedblleft{}worry\textquotedblright{}: \emph{\ru{с} \ru{тревогой}} — anxiously. \emph{\ru{Пожарная} \ru{тревога}} — a fire alarm.

\begin{cousinweb}
\emph{\ru{тревога}} is inherited from Proto-Slavic \emph{trĭvoga} (EtymDB); its deeper origin is not settled.
\end{cousinweb}

\begin{grammarlens}[title={What you've built}]
One more word for telling, explaining and describing at B1.
\end{grammarlens}

\subsection*{Your turn}
\begin{itemize}
\raggedright
  \item \emph{Say it:} \emph{trevóga}
  \item \emph{Say it:} \emph{trevóga}, once more
  \item \emph{Say it:} a sentence of your own with \emph{trevóga}
  \item \emph{Recall:} say the Russian for a stretcher, then the Russian for a siren, then say \emph{trevóga} again
\end{itemize}

\begin{tcolorbox}[breakable,colback=teal!4,colframe=teal!35!black,title={Before you move on}]
What does \ru{тревога} mean? (\textquotedblleft{}An alarm; anxiety\textquotedblright{}.) Say it once more.
\end{tcolorbox}
\section[\texorpdfstring{\ru{дежурный} (dezhúrnyy)}{дежурный (dezhúrnyy)}]{\ru{дежурный} (dezhúrnyy) — the person on duty; on duty}

\label{lesson:RU-C299-dezhurnyy}



\begin{quote}
Before the new one: say the Russian for a siren, then the Russian for an alarm.
\end{quote}

\subsection*{You'll want to know: \ru{дежурный}}
\textbf{\ru{дежурный}} (\emph{dezhúrnyy}) — \textquotedblleft{}the person on duty; on duty\textquotedblright{}.

\begin{quote}
\ru{Дежурный} \ru{врач} \ru{сразу} \ru{пришёл}.
\end{quote}

\begin{quote}
\emph{Dezhúrnyy vrach srázu prishyól.}
\end{quote}

\begin{quote}
The doctor on duty came at once.
\end{quote}

An adjective that also serves as a noun: \emph{\ru{дежурный}} — the one on duty. \emph{\ru{Дежурить}} — to be on duty.

\begin{cousinweb}
\emph{\ru{дежурный}} is built on French \emph{de jour} \textquotedblleft{}for the day\textquotedblright{} (EtymDB), with the Russian ending \emph{-\ru{ный}}: the one whose day it is.
\end{cousinweb}

\begin{grammarlens}[title={What you've built}]
One more word for telling, explaining and describing at B1.
\end{grammarlens}

\subsection*{Your turn}
\begin{itemize}
\raggedright
  \item \emph{Say it:} \emph{dezhúrnyy}
  \item \emph{Say it:} \emph{dezhúrnyy}, once more
  \item \emph{Say it:} a sentence of your own with \emph{dezhúrnyy}
  \item \emph{Recall:} say the Russian for a siren, then the Russian for an alarm, then say \emph{dezhúrnyy} again
\end{itemize}

\begin{tcolorbox}[breakable,colback=teal!4,colframe=teal!35!black,title={Before you move on}]
What does \ru{дежурный} mean? (\textquotedblleft{}The person on duty; on duty\textquotedblright{}.) Say it once more.
\end{tcolorbox}
\section[\texorpdfstring{\ru{вызывать} (vyzyvát)}{вызывать (vyzyvát)}]{\ru{вызывать} (vyzyvát) — to call (for), to send for}

\label{lesson:RU-C299-vyzyvat}



\begin{quote}
Before the new one: say the Russian for an alarm, then the Russian for the person on duty.
\end{quote}

\subsection*{You'll want to know: \ru{вызывать}}
\textbf{\ru{вызывать}} (\emph{vyzyvát}) — \textquotedblleft{}to call (for), to send for\textquotedblright{}.

\begin{quote}
\ru{Сосед} \ru{вызвал} \ru{скорую} \ru{помощь}.
\end{quote}

\begin{quote}
\emph{Soséd výzval skóruyu pómoshch.}
\end{quote}

\begin{quote}
The neighbour called an ambulance.
\end{quote}

Imperfective; perfective \emph{\ru{вызвать}} (\emph{výzvat}), \emph{\ru{вызову}}. \emph{\ru{Вызывать} \ru{врача}}, \emph{\ru{такси}}.

\begin{cousinweb}
\emph{\ru{вызывать}} is built on \emph{\ru{вы}-} \textquotedblleft{}out\textquotedblright{} and \emph{\ru{звать}} \textquotedblleft{}to call\textquotedblright{}, an inherited Slavic verb: to call someone out. EtymDB lists no parent for the compound.
\end{cousinweb}

\begin{grammarlens}[title={What you've built}]
One more word for telling, explaining and describing at B1.
\end{grammarlens}

\subsection*{Your turn}
\begin{itemize}
\raggedright
  \item \emph{Say it:} \emph{vyzyvát}
  \item \emph{Say it:} \emph{vyzyvát}, once more
  \item \emph{Say it:} a sentence of your own with \emph{vyzyvát}
  \item \emph{Recall:} say the Russian for an alarm, then the Russian for the person on duty, then say \emph{vyzyvát} again
\end{itemize}

\begin{tcolorbox}[breakable,colback=teal!4,colframe=teal!35!black,title={Before you move on}]
What does \ru{вызывать} mean? (\textquotedblleft{}To call (for), to send for\textquotedblright{}.) Say it once more.
\end{tcolorbox}
